% !TeX root = ../../../main.tex
\section{Takeuchi's Compactification of the Abel--Jacobi Morphism}\label{section:BlowupOfSymmetricPowerOfCurves}
We recall that our objective is to descend the $G$-torsor $\cP^{[d]}$ from
$U^{(d)}$ to $\PicCm[d]$ along the Abel--Jacobi map $\Phi_d$:

\[
\begin{tikzcd}
 \cP^{[d]} \arrow[d, purple]\\
 U^{(d)} \arrow[r, "\Phi_d"] & \PicCm[d]
\end{tikzcd}
\]


(Here, the purple arrow emphasizes that $\cP^{[d]}$ is a torsor on the
\'etale site of $U^{(d)}$.)

In the ramified case $\mdls>0$, \Cref{theorem:abel-jacobi-fibers}
identifies the geometric fibers of $\Phi_d$ with affine spaces. Since
$\Phi_d$ is not proper, the proper-smooth fundamental-group argument
available in the unramified case cannot be applied directly. Takeuchi
\cite[\S3]{takeuchi2018blow} overcomes this difficulty by extending $\Phi_d$
to a projective-space bundle. We now recall the ambient blowup and the
properties of the resulting compactification that will be used below.

Let $\mdls = \sum_{i=1}^s n_iP_i$, where the $P_i$ are distinct closed
points of $C$, be a modulus on $C$, and let $d$ satisfy
\Cref{equation:degree-condition}.
Takeuchi \cite[\S3, before (3.6)]{takeuchi2018blow} defines $Z_0=Z_0(\mdls,d)$ as the closed subscheme of
$C^{(d)}$ given by the image of the closed immersion
$C^{(d-\deg\mdls)}\to C^{(d)}$ that adds $\mdls$.
He also defines $X_{\mdls,d}$ as the blowup of $C^{(d)}$ along $Z_0$.
Let $E_0=E_{\mdls,d}=Z_0(\mdls,d)\times_{C^{(d)}}X_{\mdls,d}$ be the
exceptional divisor. It is irreducible of codimension~$1$, and we let
$\eta_0=\eta_{\mdls,d}$ be its generic point.


Diagrammatically:
\[\begin{tikzcd}
\overline{\{\eta_0 \}} =  E_0 \arrow[r] \arrow[d] & X_{\mdls, d} \arrow[d, "\pi"] \\
Z_0 \arrow[r, hook]           & C^{(d)}
\end{tikzcd}\]

Incorporating $U^{(d)}$, the torsor $\cP^{[d]}$, and the Abel--Jacobi map,
we have:
\[
\begin{tikzcd}
    & & \cP^{[d]} \arrow[d, purple]\\
\overline{\{\eta_0 \}} =  E_0 \arrow[r] \arrow[d] &
X_{\mdls, d} \arrow[d, "\pi"]  & \arrow[l, hook'] \arrow[dl, hook'] U^{(d)} \arrow[r, "\Phi_d"] & \PicCm[d] \\
Z_0 \arrow[r, hook]           & C^{(d)}
\end{tikzcd}
\]

We will not need Takeuchi's explicit fiber-product definition of the
compactification. The properties required below are summarized in the
following existence and structure statement, which follows from
\cite[\S3, (3.9), Proposition~3.7, and Corollary~3.8]{takeuchi2018blow}.

\begin{theorem}[Takeuchi]\label{theorem:takeuchi-compactification-theorem}
    Let $d$ satisfy \Cref{equation:degree-condition}. There exists an open
    subscheme $\Cmod{d}\subset X_{\mdls,d}$ containing $U^{(d)}$ and a
    morphism
    \[
        \widetilde\Phi_d\colon \Cmod{d}\longrightarrow\PicCm[d]
    \]
    extending $\Phi_d\colon U^{(d)}\to\PicCm[d]$, such that $\Cmod{d}$ is a
    projective space bundle over $\PicCm[d]$.
    Furthermore, the complement of $U^{(d)}$ in $\Cmod{d}$ is the scheme-theoretic intersection
    \[
        E_0 \times_{X_{\mdls,d}} \Cmod{d}
        \cong Z_0 \times_{C^{(d)}} \Cmod{d}.
    \]
\end{theorem}



Diagrammatically we have:
\[
\begin{tikzcd}
    &  & \cP^{[d]} \arrow[d, purple]\\
    E_0 \times_{X_{\mdls,d}} \Cmod{d} \arrow[d] \arrow[r]  & \Cmod{d} \arrow[d, hook] \arrow[rd, "\widetilde\Phi_d"]  & U^{(d)} \arrow[l, hook'] \arrow[d, "\Phi_d"] \\
\overline{\{\eta_0 \}} =  E_0 \arrow[r] \arrow[d] &
X_{\mdls, d} \arrow[d, "\pi"]  &  \PicCm[d] &  \\
Z_0 \arrow[r, hook]           & C^{(d)}
\end{tikzcd}
\]

Takeuchi's description of the compactification identifies its boundary as a
relative hyperplane; see
\cite[Lemmas~3.2 and~3.3(1), Theorem~3.4, (3.9), and
Corollary~3.8]{takeuchi2018blow}. This gives the following fiberwise
description.

\begin{cor}[The boundary is a relative hyperplane]
\label{lemma:geometric-generic-fiber-takeuchi}
The boundary
\[
    H=\Cmod{d}\setminus U^{(d)}
\]
is a smooth relative hyperplane divisor in the projective-space bundle
$\widetilde\Phi_d:\Cmod{d}\to\PicCm[d]$. Consequently, if
$\bar s=\Spec\Omega\to\PicCm[d]$ is any geometric point and
$r=d-\deg\mdls-g+1$, then there is an isomorphism of pairs\footnote{For pairs
$(X,D)$ and $(Y,E)$, an isomorphism of pairs is an isomorphism
$f\colon X\xrightarrow{\sim}Y$ such that
$D=X\times_Y E$ as closed subschemes of $X$. No canonical choice of $f$ is
asserted here.}
\[
    \left((\Cmod{d})_{\bar s},H_{\bar s}\right)
    \cong
    \left(\mathbb P^r_{\Omega},\mathbb P^{r-1}_{\Omega}\right),
\]
and hence $U^{(d)}_{\bar s}\cong\mathbb A^r_{\Omega}$.
\end{cor}

\begin{proof}
The formation of a projective bundle and of its relative hyperplane commutes
with arbitrary base change. Hence, after base change to $\bar s$, the
projective-space bundle $\Cmod{d}$ becomes $\mathbb P^r_{\Omega}$ and $H$
becomes a hyperplane $\mathbb P^{r-1}_{\Omega}$. Its complement is therefore
$\mathbb A^r_{\Omega}$.
\end{proof}

